\documentclass[11pt]{article}

% \newcommand{\answerDirectory}{solutions}

../../macros/macros.tex

\inputAnswer{problem_0_answer}

\doHwHeader{Homework \BLUE{4}}

%%%%% EDIT THE FILES NAMED problem_X_answer.tex.  DON'T EDIT THIS FILE

%%%%%%%%%%%%%%%%%%%%%%%%%%%%%%%%%%%%%%%%%%%%%%%%%%%%%%%%%%%% 
%%%%%%%%%%%%%%%%%%%%%%%%%%%%%%%%%%%%%%%%%%%%%%%%%%%%%%%%%%%% 
%%%%%%%%%%%%%%%%%%%%%%%%%%%%%%%%%%%%%%%%%%%%%%%%%%%%%%%%%%%% 
%%%%%%%%%%%%%%%%%%%%%%%%%%%%%%%%%%%%%%%%%%%%%%%%%%%%%%%%%%%% 

\begin{document}


%%%%% EDIT THE FILES NAMED problem_X_answer.tex.  DON'T EDIT THIS FILE 

\paragraph{General instructions.}
% First read \BLUE{Lecture Note 4}, then
Answer the problems given in the rest of this document.
Edit this LaTeX template to add your answers, as described in the file named \texttt{00\_README.txt},
then compile it to produce \BLUE{\texttt{homework\_4.pdf}}\ifGradescope{} for submission to gradescope\fi.

\ifGradescope
For this homework, when you submit to gradescope you don't need to identify the pages
for each problem.  We're using gradescope's ``fixed-length'' grading mode.
To make it work, your answer for each part should fit in the allotted space,
so that each subsequent answer is on the expected page in the expected location.
Your PDF should have
9
pages.  
\fi

\paragraph{Don't know?  Say so.}
If you don't know the answer to a problem or any part of the problem,
we encourage you to simply write \emph{``I don't know''} instead of giving an answer.
(If you write ``I don't know'' you can also add any questions you might have.)

\paragraph{Gentle reminder: homework is for learning.}
As noted in \LN{prelim}, engaging fully with the homeworks is the best way to learn the material.
\emph{The primary role of homeworks is to help you learn the material,
  and the primary role of feedback on homeworks is to help you improve your ideas.}

Please try to explain your ideas as clearly as you can.
The clearer your explanations are, the more useful the feedback on them can be.

\paragraph{Collaboration and citation policy.}
If you are doing this homework as part of a course, follow the course's collaboration and citation policy.
In any case, at the end of each answer, list the people you collaborated with,
and cite every source that your answer uses ideas from
(other than the lecture notes and the sources listed in the lecture notes).
If there are none, write ``none''.

%%%%% EDIT THE FILES NAMED problem_X_answer.tex.  DON'T EDIT THIS FILE 

\vfill
\noindent{\footnotesize\copyrightNotice}

%%% Local Variables:
%%% mode: latex
%%% TeX-master: "homework_2"
%%% End:


%%%%%%%%%%%%%%%%%%%%%%%%%%%%%%%%%%%%%%%%%%%%%%%%%%%%%%%%%%%% 
%%%%%%%%%%%%%%%%%%%%%%%%%%%%%%%%%%%%%%%%%%%%%%%%%%%%%%%%%%%% 
%%%%%%%%%%%%%%%%%%%%%%%%%%%%%%%%%%%%%%%%%%%%%%%%%%%%%%%%%%%% 
%%%%%%%%%%%%%%%%%%%%%%%%%%%%%%%%%%%%%%%%%%%%%%%%%%%%%%%%%%%% 

%%%%% EDIT THE FILES NAMED problem_X_answer.tex.  DON'T EDIT THIS FILE 

%%%%%%%%%%%%%%%% PROBLEM 1  %%%%%%%%%%%%%%%%%

% \setcounter{problem}{1}

\newpage 
\problem \emph{(Section 5.3)}
Consider executing the \prob{Activity Selection} algorithm \textsf{rselect} (as defined in Section 5.3)
on the following instance of the \prob{Activity Selection} problem:

\[I = \{[0, 4], [2, 6], [3, 8], [7, 11], [9, 13], [10, 14], [12, 15], [12, 16], [14, 18]\}.\]

\centerline{
  \begin{tikzpicture}[scale=0.5]\footnotesize
    \draw (0,0) -- (4, 0) node[midway,auto] {\footnotesize $J_1$}; 
    \draw (2,1) -- (6, 1) node[midway,auto] {\footnotesize $J_{2}$};
    \draw (3,2) -- (8, 2) node[midway,auto] {\footnotesize $J_{3}$};
    \draw (7,0) -- (11, 0) node[midway,auto] {\footnotesize $J_{4}$};
    \draw (9, 1) -- (13, 1) node[midway,auto] {\footnotesize $J_{5}$}; 
    \draw (10,2) -- (14, 2) node[midway,auto] {\footnotesize $J_{6}$};
    \draw (12,3) -- (15, 3) node[midway,auto] {\footnotesize $J_{7}$};
    \draw (12,0) -- (16 , 0) node[midway,auto] {\footnotesize $J_{8}$};
    \draw (14,1) -- (18 , 1) node[midway,auto] {\footnotesize $J_{9}$};
  \end{tikzpicture}
}

\begin{enumerate}[label=(\alph*)]

\item When \textsf{rselect}$(I)$ executes, what jobs are in the set $D$ that it computes in Line 3,
  before recursing?

\item What set of jobs does \textsf{rselect}$(I)$ output?

\item Consider running the linear-time iterative implementation, \textsf{iselect} (Section 5.3, second definition)
  on input $I$.  What values does the variable $\ell$ take on during the execution?  Give them in increasing order.

\item Consider the following ``theorem'' and ``proof''.  Is the proof correct?
  If not, what is the first step in the proof that makes an assertion that does not necessarily hold?

  \begin{theorem}
    Consider any instance $I'=\{[s_1, e_1],  [s_2, e_2], \ldots, [s_n, e_n]\}$ of \prob{Activity Selection}.
    Let $X=\{x_1, x_2, \ldots, x_k\}$ be a set of times such that,
    for every job $[s_i, e_i]$ in $I'$, the interval $[s_i, e_i]$ contains at least one of the times in $X$.
    Then no pairwise-disjoint subset of $I'$ contains more than $k$ jobs.
  \end{theorem}
  \begin{longFormProof}
    \step Consider any such $I'=\{[s_1, e_1],  [s_2, e_2], \ldots, [s_n, e_n]\}$ and $X=\{x_1, x_2, \ldots, x_k\}$.
    \begin{block}\label{block: as contra}
      \step Assume for contradiction that some pairwise-disjoint subset $S$ of $I'$ has more than $k$ jobs.
      \step By the assumption on $X$ in the theorem, each job in $S$ contains at least one point in $X$.
      \step There are more than $k$ jobs in $S$, but only $k$ points in $X$.
      \step So (by the previous two steps) some point $x_i$ in $X$ is in at least two jobs in $S$.
      \step So there are two jobs in $S$ that are not disjoint (as they each contain $x_i$).
      \step This contradicts that $S$ is pairwise-disjoint.
    \end{block}
    \step By Block~\ref{block: as contra}, no pairwise-disjoint subset of $I'$ contains more than $k$ jobs.
  \end{longFormProof}

\item Is the following conjecture true or false?  Don't justify your answer.

  \begin{conjecture}
    Consider any instance $I'=\{[s_1, e_1],  [s_2, e_2], \ldots, [s_n, e_n]\}$ of \prob{Activity Selection}.
    Let $e'_1, e'_2, \ldots, e'_k$ be the end times of the jobs that are chosen by the algorithm.
    Then, for every job $[s_i, e_i]$ in $I'$, the interval $[s_i, e_i]$ contains at least one of the end times
    $e'_1, e'_2, \ldots, e'_k$.
  \end{conjecture}  

  (Note that if the theorem and conjecture are both true,
  together they imply that no pairwise-disjoint subset of $I'$
  is larger than the one returned by the greedy algorithm.)
  
\end{enumerate}

\newpage

\setlength{\ITEMHT}{1.5in}
{
  \injectEnumerate 

  \inputAnswer{problem_1_answer}
}

%%%%%%%%%%%%%%%%%%%%%%%%%%%%%%%%%%%%%%%%%%%%%%%%%%%%%%%%%%%% 
%%%%%%%%%%%%%%%%%%%%%%%%%%%%%%%%%%%%%%%%%%%%%%%%%%%%%%%%%%%% 
%%%%%%%%%%%%%%%%%%%%%%%%%%%%%%%%%%%%%%%%%%%%%%%%%%%%%%%%%%%% 
%%%%%%%%%%%%%%%%%%%%%%%%%%%%%%%%%%%%%%%%%%%%%%%%%%%%%%%%%%%% 

%%%%% EDIT THE FILES NAMED problem_X_answer.tex.  DON'T EDIT THIS FILE 

%%%%%%%%%%%%%%%% PROBLEM 2  %%%%%%%%%%%%%%%%%

\newpage
\problem\emph{(Section 5.2)}
{\small
  For each $d\in\{6,7,8,9,11\}$, 
  consider \prob{Making Change} with denominations (coin values) $\{1, 5, d\}$.
  For one of these five values, 
  it is true that, for all integers $N\ge 0$, the greedy algorithm (using available denominations $\{1, 5, d\}$)
  makes change for $N$ using a minimum possible number of coins.
  For the other four values, it is not true.

  (a) Fill out the table in the template showing, for each value of $d$,
  the smallest $N$ for which the greedy algorithm uses too many coins.
  (For the $d$ for which the greedy algorithm always uses a minimum number of coins, put $\infty$.)

  (b) For the value of $d$ for which the greedy algorithm can't use too many coins,
  give a proof of the following lemma for that particular value of $d$:

  \begin{lemma}
    For any $N\ge d$, \emph{some} optimal way of making change for $N$ (using coin values in $\{1, 5, d\}$)
    contains at least one $d$.
  \end{lemma}

  Note the emphasis on ``some''.
}
\bigskip

\lineacross

% \continued

\setlength{\ITEMHT}{2in}
{
  % \injectEnumerate 

  \inputAnswer{problem_2_answer}
}

\newpage

%%%%%%%%%%%%%%%%%%%%%%%%%%%%%%%%%%%%%%%%%%%%%%%%%%%%%%%%%%%% 
%%%%%%%%%%%%%%%%%%%%%%%%%%%%%%%%%%%%%%%%%%%%%%%%%%%%%%%%%%%% 
%%%%%%%%%%%%%%%%%%%%%%%%%%%%%%%%%%%%%%%%%%%%%%%%%%%%%%%%%%%% 
%%%%%%%%%%%%%%%%%%%%%%%%%%%%%%%%%%%%%%%%%%%%%%%%%%%%%%%%%%%% 

%%%%% EDIT THE FILES NAMED problem_X_answer.tex.  DON'T EDIT THIS FILE 

%%%%%%%%%%%%%%%% PROBLEM 3  %%%%%%%%%%%%%%%%%

\newpage
\problem \emph{(see Exercise 5.14)}\emph{(Greedy step for \prob{Puncturing Intervals})} {\small

  The \prob{Puncturing Intervals} problem is defined as follows: 

  \begin{mylist}
  \item[\bf input:] Set $I=\{J_1, J_2, \ldots, J_n\}$ of intervals, where $J_i = [s_i, e_i]$.
    We assume $e_1 \le e_2 \le \cdots \le e_n$.
  \item[\bf output:] A minimum-size finite set $P=\{p_1, p_2, \ldots, p_k\}$
    of points such that $J_i \cap P \ne \emptyset$ for each $J_i\in I$.
  \end{mylist}

  That is, you want to choose a minimum-size set of points
  such that every given interval $J_i$ contains at least one of those points.
  For example, for the set $I$ of intervals in Problem 1,
  the set $P=\{1, 5, 9, 14\}$ of points does have the property that
  it contains at least one point in each interval in $I$.
  However, that set $P$ is not a correct solution, because it doesn't have minimum possible size.

  \begin{enumerate}[label=(\alph*)]
  \item
    Give a correct solution $P$ for the % \prob{Minimum Interval Pinner}
    instance defined by that set $I$ of intervals.

  \item
    Describe a rule that, given an instance $I$ with $e_1 \le e_2\le \cdots \le e_n$,
    chooses a single point $p_1$ such that
    there is guaranteed to exist a correct solution $P$ (for the instance $I$) that contains $p_1$.
    Describe your rule precisely in a single sentence.
    Your rule should be implementable as an algorithm that takes $I$
    as input and produces the desired point $p_1$ in \emph{constant} time.
    Do not explain your reasoning for why you think your rule works.

  \end{enumerate}
}

\lineacross

\setlength{\ITEMHT}{1.5in}
{
  \injectEnumerate 

  \inputAnswer{problem_3_answer}
}

%%%%%%%%%%%%%%%%%%%%%%%%%%%%%%%%%%%%%%%%%%%%%%%%%%%%%%%%%%%% 
%%%%%%%%%%%%%%%%%%%%%%%%%%%%%%%%%%%%%%%%%%%%%%%%%%%%%%%%%%%% 
%%%%%%%%%%%%%%%%%%%%%%%%%%%%%%%%%%%%%%%%%%%%%%%%%%%%%%%%%%%% 
%%%%%%%%%%%%%%%%%%%%%%%%%%%%%%%%%%%%%%%%%%%%%%%%%%%%%%%%%%%% 

%%%%% EDIT THE FILES NAMED problem_X_answer.tex.  DON'T EDIT THIS FILE 

\end{document}

%%% Local Variables:
%%% mode: latex
%%% TeX-master: t
%%% End:
